\documentclass[11pt,letterpaper]{article}

\usepackage[margin=2.5cm]{geometry}
\usepackage{amsmath}
\usepackage{amssymb}
\usepackage{bm}
\usepackage{booktabs}
\usepackage{array}
\usepackage{ragged2e}
\usepackage{seqsplit}
\usepackage{enumitem}
\usepackage{xcolor}
\usepackage{listings}
\usepackage[hidelinks]{hyperref}
\usepackage{caption}
\usepackage{float}
\usepackage{graphicx}

\DeclareGraphicsExtensions{.pdf}
\graphicspath{{figs/}}

\DeclareRobustCommand{\file}[1]{\texttt{\seqsplit{#1}}}
\newcommand{\code}[1]{\texttt{#1}}
\setlength{\emergencystretch}{3em}
\newcolumntype{L}[1]{>{\RaggedRight\arraybackslash}p{#1}}

% Symbols used throughout.  Bold lower case = a column of n numbers, one per
% channel; bold upper case = a matrix.
\newcommand{\va}{\mathbf{a}}
\newcommand{\vx}{\mathbf{x}}
\newcommand{\vr}{\mathbf{r}}
\newcommand{\vd}{\mathbf{d}}
\newcommand{\vw}{\mathbf{w}}
\newcommand{\vz}{\mathbf{z}}
\newcommand{\vs}{\mathbf{s}}
\newcommand{\vv}{\mathbf{v}}
\newcommand{\vone}{\mathbf{1}}
\newcommand{\vu}{\mathbf{u}}
\newcommand{\epsy}{\bm{\varepsilon}_y}
\newcommand{\epsb}{\bm{\varepsilon}_b}
\newcommand{\ey}{\mathbf{e}_y}
\newcommand{\eb}{\mathbf{e}_b}
\newcommand{\mA}{\mathbf{A}}
\newcommand{\mX}{\mathbf{X}}
\newcommand{\mP}{\mathbf{P}}
\newcommand{\mI}{\mathbf{I}}
\newcommand{\mE}{\mathbf{E}}
\newcommand{\mF}{\mathbf{F}}
\newcommand{\mM}{\mathbf{M}}
\newcommand{\mL}{\mathbf{L}}
\newcommand{\mC}{\mathbf{C}}
\newcommand{\mB}{\mathbf{B}}
\newcommand{\mU}{\mathbf{U}}
\newcommand{\mS}{\mathbf{S}}
\newcommand{\mV}{\mathbf{V}}
\newcommand{\T}{^{\mathsf{T}}}
\newcommand{\inv}{^{-1}}

% Code listings.  Line numbers are the real line numbers in duotone.py as of
% the date on the title page.
\definecolor{codebg}{gray}{0.96}
\lstdefinestyle{py}{
  language=Python,
  basicstyle=\ttfamily\footnotesize,
  backgroundcolor=\color{codebg},
  frame=single,
  numbers=left,
  numberstyle=\tiny\color{gray},
  breaklines=true,
  breakatwhitespace=true,
  showstringspaces=false,
  keywordstyle=\color{blue!55!black},
  commentstyle=\color{green!35!black},
  stringstyle=\color{red!45!black},
  captionpos=b,
}
\newcommand{\src}[3]{%
  \lstinputlisting[style=py,firstline=#1,lastline=#2,firstnumber=#1,
     caption={#3}]{./figs/duotone-v2026-09-27.py}}
%    caption={#3}]{../desktop/padanalysis/duotone.py}}

\title{\file{duotone.py}: splitting a pad into yellow and brown\\
\large What it does, why, and the matrix algebra behind it, with a review}
\author{Colorimeter project --- pad analysis}
\date{2026-09-28, branch \file{preimage-capture}}

\begin{document}
\maketitle

\begin{abstract}
\noindent
\file{desktop/padanalysis/duotone.py} takes the photographs of one test pad
and, for every pixel, estimates how much of a yellow color and how much of a
brown color it holds.  This report explains, step by step and in plain
language, what each function computes and why, and writes out the matrix
algebra so each step can be followed.  For each matrix product it says what
the product \emph{means}.  All the figures and numbers come from real pads
(C1017--C1025 and C1028--C1036) processed by the module as it stands today.

The report is also a review.  Section~\ref{sec:findings} lists what is wrong
or misleading.  The most important finding: the step that finds the two
reference colors from the pixels, \code{endmembers()}, currently returns
colors almost exactly \emph{opposite} to the real yellow and brown
(Figure~\ref{fig:spectra}).  Everything after it inherits that error.  The
remaining steps are sound, with smaller issues noted.  This report does not
change any code; the fixes are proposals for you to decide on.
\end{abstract}

\tableofcontents
\newpage

%======================================================================
\section{What the module is for}
%======================================================================

The Gutzeit reaction leaves a stain on the mercuric bromide pad.  At low
arsenic the stain is pale yellow; at high arsenic the center of the dot turns
brown.  These are two different chemical products with two different colors,
not one color at two strengths.  One number per pad cannot describe both, so
\file{duotone.py} works pixel by pixel and asks, for each pixel: \emph{how
much yellow, and how much brown?}

The module is a chain of small functions.  Table~\ref{tab:functions} lists
them in the order the data flows through them, and who calls each one today.

\begin{table}[H]
\centering\small
\begin{tabular}{@{}L{3.3cm}L{7.2cm}L{4.3cm}@{}}
\toprule
Function & What it does & Used by (today) \\
\midrule
\code{PadCube} & Loads one pad: absorbance for every pixel in every
  channel, plus a noise level per channel. (Sec.~\ref{sec:data}) &
  \code{cell\_PCA}, \code{pca\_study}, \code{build} \\
\code{contrast} & Removes from each pixel the part that is the same in every
  channel (shading). (Sec.~\ref{sec:contrast}) &
  \code{cell\_PCA}, \code{pca\_study}, \code{fit}, \code{unmix} \\
\code{hue\_plane} & Finds the flat sheet (plane) in which nearly all pixels
  lie; this is the ``color map''. (Sec.~\ref{sec:plane}) &
  \code{cell\_PCA}, \code{pca\_study}, \code{fit} \\
\code{endmembers} & Picks the two reference colors, yellow and brown, as the
  two edges of the cloud of pixels on that map. (Sec.~\ref{sec:endmembers}) &
  \code{cell\_PCA} (by default), \code{pca\_study}, \code{fit} \\
\code{unmix} & Splits each pixel into a yellow amount and a brown amount.
  (Sec.~\ref{sec:unmix}) & \code{cell\_PCA}, \code{\_\_main\_\_} \\
\code{residual\_rms} & How much of the pixels the two colors fail to explain.
  & \code{cell\_PCA}, \code{\_\_main\_\_} \\
\code{classify} & Labels each pixel white, yellow or brown.
  (Sec.~\ref{sec:classify}) & \code{\_\_main\_\_} \\
\code{amount\_sigma}, \code{noise\_cov} & How noisy the yellow and brown
  amounts are. (Sec.~\ref{sec:noise}) & \code{\_\_main\_\_} \\
\code{j\_hue}, \code{j\_amounts}, \code{subset\_sigma},
  \code{subset\_unmix} & Score a chosen subset of channels: how well it
  separates or measures the two colors. (Sec.~\ref{sec:subsets}) &
  \code{report14} only (a record; no longer runs) \\
\code{build}, \code{fit} & Load many pads on a shared channel list; find one
  pair of reference colors for all of them. (Sec.~\ref{sec:helpers}) &
  \code{\_\_main\_\_}, \code{look3} (record) \\
\code{render\_rgb} & A true-color picture, for looking at, not measuring.
  (Sec.~\ref{sec:helpers}) & \code{cell\_PCA} \\
\bottomrule
\end{tabular}
\caption{The functions of \file{duotone.py}, in data-flow order.
\code{cell\_PCA} is in \file{analysis\_1.py}; \code{pca\_study} is
\file{pca\_study.py}; \code{\_\_main\_\_} is what runs when
\file{duotone.py} is run by itself.}
\label{tab:functions}
\end{table}

\paragraph{A note on the name ``PCA''.}  The code and printouts speak of
``PCA'' and ``PC1, PC2''.  Strictly, what it does is a close relative of
principal component analysis (PCA) that measures spread from the origin
rather than from the average pixel.  Section~\ref{sec:plane} explains the
difference and why the relative is the right choice here.

%======================================================================
\section{The data: one pixel is a list of six numbers}
\label{sec:data}
%======================================================================

\subsection{Channels}

A \emph{channel} (the code says \emph{cell}) is one LED seen through one
color layer of the camera's Bayer filter, for example
\code{royal\_blue.B}: the royal-blue LED, read from the blue-filtered
pixels.  Each channel reports how much light of roughly one color the pad
reflects.

Only channels that the pre-image lit can be used, because only those have a
white reference (Section~\ref{sec:absorbance}).  All 18 pads share six:
\begin{center}\small
\begin{tabular}{@{}lllllll@{}}
\toprule
& violet.B & royal\_blue.B & cyan.B & green.Gb & green.Gr & photo\_red.R \\
\midrule
LED peak (nm) & 410 & 455 & 500 & 525 & 525 & 660 \\
\bottomrule
\end{tabular}
\end{center}
Throughout this report $n$ is the number of channels ($n = 6$ here) and $N$
is the number of pad pixels (about 2{,}200 per pad).  The two green channels
are the same LED seen through the camera's two green filter layers; they are
kept apart because they differ slightly.

\subsection{From camera counts to absorbance}
\label{sec:absorbance}

\code{PadCube} does not compute anything new here; it calls
\code{analysis\_1.get\_balanced\_cell\_data}, which does this for every
channel:
\begin{enumerate}[nosep]
\item subtract the camera's dark level;
\item divide by the gray chips (so the lamp's brightness cancels):
  this gives the reflectance $R$;
\item divide by the median $R$ of the same strip's \emph{dry pad in the
  pre-image} (the photo taken before any reagent): this gives
  $R_\infty = R / R_\text{white}$, which is 1 on an unstained pad;
\item take $A = -\log_{10} R_\infty$, the \emph{absorbance}.
\end{enumerate}
So $A = 0$ means ``as white as the dry pad was'', and $A = 0.3$ means ``this
pixel returns half as much light as it did before the test''
($10^{-0.3} \approx 0.5$).

\paragraph{Why absorbance, not reflectance?}  Two reasons, both about
\emph{adding}.
\begin{itemize}[nosep]
\item \textbf{Shading adds a constant.}  A shadow, or a pixel that sits a
  little lower in the paper, returns a fraction $s$ of the light in
  \emph{every} channel: $R_\infty \to s R_\infty$.  In absorbance that is
  $A \to A - \log_{10} s$: the \emph{same number added to every channel}.
\item \textbf{Dyes add (approximately).}  If two colored substances both
  absorb, their absorbances add (Beer--Lambert).  This is exact for light
  passing through a clear solution; for light scattered back out of paper it
  is only approximate (Section~\ref{sec:limits}).
\end{itemize}
Both effects then appear as sums, and sums are what matrix algebra handles.

\subsection{The matrix $\mA$}

\code{PadCube.A} holds the absorbances as a table with one row per pixel and
one column per channel($N$ pixels, $n$ channels):
\[
\mA = \begin{pmatrix}
A_{1,1} & A_{1,2} & \cdots & A_{1,n}\\
\vdots  &         &        & \vdots \\
A_{N,1} & A_{N,2} & \cdots & A_{N,n}
\end{pmatrix}
\quad (N \times n).
\]
Row $i$ is pixel $i$'s \emph{spectrum}: six absorbances, one per channel.
When the report talks about one pixel it writes that row as a column
$\va = (a_1, \dots, a_n)\T$.  The rest of \code{PadCube} is bookkeeping:
\code{mask} says which pixels of the image are the rows, \code{yx} gives
each row's position, and \code{paint()} writes a per-pixel result back into
an image.  Pixels are left out if they are known hot pixels or if any
channel lacks a value there, so every channel has the same rows.

\subsection{Pixel scatter per channel}
\label{sec:padnoise}

To know whether a small absorbance is real, \code{PadCube} measures how much
each channel's pixels scatter from one pixel to the next on the dry pad in
the pre-image, where there is no stain.  The code calls this scatter
``noise''.  Section~\ref{sec:scatterwhat} shows it is mostly the paper's own
texture, not camera noise.

\src{172}{179}{\code{PadCube}: each channel's dry pad, read from the
pre-image on its own grid, and turned into a plain array.}

\src{196}{202}{\code{PadCube}: pixel scatter from neighbor differences.}

For each pair of side-by-side pixels it takes the difference and divides by
$\sqrt2$:
\[
\delta = \frac{a_{\text{right}} - a_{\text{left}}}{\sqrt 2}.
\]
Why a difference: the pad has slow, smooth changes across it (light
fall-off, the paper's own gradient).  Two neighbors share almost the same
smooth part, so it cancels in their difference, and what is left is the
pixel-to-pixel scatter.  Why $\sqrt2$: if each pixel scatters independently
by $\sigma$, the difference of two scatters by $\sigma\sqrt2$; dividing by
$\sqrt2$ gives back $\sigma$.

\src{252}{256}{\code{\_pair\_noise}: a robust spread.}

The spread is measured with the \emph{median absolute deviation} (MAD): the
median of $|\delta - \text{median}(\delta)|$.  One speck of dust changes a
median very little, where it could change a standard deviation a lot.  For
ordinary bell-curve scatter, $\sigma = 1.4826 \times \text{MAD}$, hence the
constant.  On these pads $\sigma$ is 0.006--0.009 in absorbance for every
channel.  \code{PadCube.dnoise} keeps all the differences so that
\code{noise\_cov} can later see whether the channels scatter together
(Section~\ref{sec:noise}).

\subsubsection{Reading the numpy}

Two pieces of notation unlock these lines.

\paragraph{NaN.}  Each channel's pad arrives from \code{analysis\_1} as a
\emph{masked array}: the pixels off the pad carry a ``no value'' flag.
\code{np.ma.filled(..., np.nan)} swaps each flag for NaN (``not a number''),
so from then on everything is a plain array and NaN means ``not on the pad''.

\paragraph{Slicing.}  \code{a[:, 1:]} means ``all rows (\code{:}),
columns from 1 to the end''; \code{a[:, :-1]} means ``all rows, columns from
0 up to but not including the last''.  For one row of four pixels:
\begin{center}\small
\begin{tabular}{@{}ll@{}}
\code{a[:, :-1]} & $(p_0,\ p_1,\ p_2)$ \quad each pixel that has a right-hand neighbor \\
\code{a[:, 1:]}  & $(p_1,\ p_2,\ p_3)$ \quad that neighbor \\
difference       & $(p_1 - p_0,\ p_2 - p_1,\ p_3 - p_2)$
\end{tabular}
\end{center}
The two slices line up so that position $j$ holds a pixel and its right-hand
neighbor.  Table~\ref{tab:numpy} goes through the lines one by one.

\begin{table}[tbp]
\centering\small
\begin{tabular}{@{}l L{4.6cm} L{6.6cm} L{2.0cm}@{}}
\toprule
Line & Code & What it does & Result shape \\
\midrule
178--179 & \code{np.ma.filled(pre.A\allowbreak.astype(float), np.nan)} &
  \code{astype(float)} makes the values floating point; \code{filled}
  replaces every ``no value'' flag with NaN. & (rows, cols) \\
196 & \code{np.stack(pre\_planes, axis=-1)} & Piles the $n$ channel images
  into one 3-D block.  \code{axis=-1} means ``put the new axis last'', so
  \code{pre[i, j, k]} is pixel $(i, j)$ in channel $k$. & (rows, cols, $n$) \\
197 & \code{np.isfinite(pre)\allowbreak.all(axis=-1)} & \code{isfinite} is
  True where there is a number, False at NaN.  \code{.all(axis=-1)} combines
  across the channels: True only if \emph{every} channel has a value there. &
  (rows, cols) \\
198 & \code{ok[:, :-1] \& ok[:, 1:]} & The slicing above.  \code{\&} is
  ``and'', element by element: True where a pixel \emph{and} its right-hand
  neighbor are both on the pad. & (rows, cols$-1$) \\
199 & \code{pre[:, 1:, :] -\allowbreak{} pre[:, :-1, :]} & Right minus left,
  for every pair and every channel.  The third \code{:} keeps all channels. &
  (rows, cols$-1$, $n$) \\
199 & \code{[pair]} & Indexing with a True/False array shaped like the first
  two axes: keeps only the True positions and lists them one after another,
  row by row.  The channel axis is kept. & ($n_\text{pairs}$, $n$) \\
200 & \code{/ np.sqrt(2.0)} & Divides every entry by $\sqrt2$. & same \\
201--202 & \code{np.array([\_pair\_noise(\allowbreak self.dnoise[:, k])
  for k in range(len(cells))])} & A \emph{list comprehension}: for each
  channel $k$, take column $k$ (all pairs, that channel), reduce it to one
  number, and collect the $n$ numbers into an array. & ($n$,) \\
252 & \code{d.size} & The number of entries; if there are none, the answer is
  NaN. & \\
255 & \code{np.median(np.abs(\allowbreak d - np.median(d)))} & From the
  inside out: the median of \code{d}; subtract it from every entry; take
  absolute values; take the median of those.  This is the MAD. & one number \\
256 & \code{float(1.4826 * mad)} & Converts the MAD to $\sigma$.
  \code{float()} turns numpy's number type into a plain Python number. & one
  number \\
\bottomrule
\end{tabular}
\caption{The numpy in the scatter measurement, line by line.  Line numbers
are those of \file{duotone.py}.}
\label{tab:numpy}
\end{table}

\subsubsection{Is this the right way to measure it?}
\label{sec:scatterwhat}

\paragraph{The method is standard.}  Estimating scatter from differences
between neighbors is a textbook method: Rice introduced the one-dimensional
version for statistics,\footnote{J.~Rice, ``Bandwidth choice for
nonparametric regression'', \emph{Annals of Statistics} 12 (1984)
1215--1230.} and for images Immerk{\ae}r uses a $3\times3$ kernel close to
``each pixel minus the mean of its four neighbors''.\footnote{J.~Immerk{\ae}r,
``Fast noise variance estimation'', \emph{Computer Vision and Image
Understanding} 64 (1996) 300--302.}  Both are unbiased when each pixel
scatters independently of its neighbors.  (The four-neighbor version
divides by $\sqrt{1.25}$ instead of $\sqrt2$, because the neighbors' average
itself scatters by $\sigma/2$.)  A smooth gradient across the pad adds the
same constant to every difference, and the MAD's subtraction of the median
removes it.

\paragraph{Other directions and four neighbors give the same answer.}  On
four dry pads (C1017, C1025, C1032, C1034), all six channels, relative to the
horizontal pairs used here:
\begin{center}\small
\begin{tabular}{@{}lc@{}}
\toprule
method & $\sigma$ relative to horizontal pairs (average of 24) \\
\midrule
vertical pairs & 1.015 \\
diagonal pairs & 1.057 \\
pixel minus the mean of its 4 neighbors & 0.970 \\
\bottomrule
\end{tabular}
\end{center}
Each estimate uses about 1{,}200 pairs, so it is itself uncertain by about
$\pm4\,\%$.  The methods agree within that.

\paragraph{What it measures is mostly the paper's texture.}  The shading
study (\file{docs/shading\_study\_2026-09-29.md}, step~1a) found that 90--99\,\%
of the pixel-to-pixel variation repeats from one exposure to the next, and
that it moves with the strip when the strip is moved.  Camera noise is only
about 0.002 in absorbance.  So $\sigma$ is the grain of the paper.  That is
still the right yardstick for ``does this pixel hold dye'': a pixel has to
stand out from the paper's own grain, not just from the camera's noise.  It
also means the pre-image's shorter exposure hardly matters.

\paragraph{Neighbors are not independent.}  The texture is 1--2~px across:
neighboring pixels on the dry pad correlate about 0.2 at 1~px and 0.1 at
2~px.  Whatever two neighbors share cancels in their difference, so
$\sigma$ may read up to about 10\,\% low as a measure of the full scatter
(all neighbor methods share this).  It also means an average over many
pixels shrinks more slowly than for independent pixels.

%======================================================================
\section{The model}
\label{sec:model}
%======================================================================

The module assumes each pixel's spectrum is built from three ingredients:
\begin{equation}
\va \;=\; g\,\vone \;+\; y\,\epsy \;+\; b\,\epsb \;+\; \text{noise}.
\label{eq:model}
\end{equation}
\begin{itemize}[nosep]
\item $\vone = (1,1,\dots,1)\T$: the same amount in every channel.  Its
  amount $g$ is the \emph{shading}: any shadow or other effect that dims all
  channels equally.  $g$ may be positive or negative, and is thrown away.
\item $\epsy$: the yellow spectrum, how strongly the yellow product absorbs
  in each channel.  Its amount $y \ge 0$.
\item $\epsb$: the brown spectrum.  Its amount $b \ge 0$.
\end{itemize}
Amounts cannot be negative: there is no such thing as a negative amount of
dye.  The two spectra are the same for every pixel of every pad; only the
three amounts change from pixel to pixel.

\paragraph{Matrix form.}  Put the three ingredients side by side as the
columns of one matrix, and the three amounts in one column:
\begin{equation}
\va = \mF\,\bm\theta,\qquad
\mF = \begin{pmatrix} \vone & \epsy & \epsb \end{pmatrix}\ (n \times 3),
\qquad
\bm\theta = \begin{pmatrix} g \\ y \\ b \end{pmatrix}.
\label{eq:modelmatrix}
\end{equation}
\emph{Meaning of the product $\mF\bm\theta$:} $\mF$ is the list of
ingredients, $\bm\theta$ is the recipe (how much of each), and the product
is the spectrum that recipe would produce.  Fitting the model means finding,
for each pixel, the recipe whose product comes closest to the measured
$\va$.

%======================================================================
\section{Step 1: remove the part that is the same in every channel}
\label{sec:contrast}
%======================================================================

\src{259}{272}{\code{flat\_direction} and \code{contrast}.}

\subsection{The algebra}

Let $\vu = \vone/\sqrt n$, the all-ones vector scaled to length 1 (this is
\code{flat\_direction}).  For one pixel:
\begin{itemize}
\item $\vu\T\va = \frac{1}{\sqrt n}\sum_k a_k = \sqrt n\,\bar a$, where
  $\bar a$ is the pixel's average over its channels.  \emph{Meaning:} how far
  the pixel reaches in the ``same in every channel'' direction.
\item $(\vu\T\va)\,\vu = \sqrt n\,\bar a \cdot \frac{1}{\sqrt n}\vone =
  \bar a\,\vone$.  \emph{Meaning:} the flat part of the pixel, a spectrum with
  $\bar a$ in every channel.
\item The contrast is what is left:
  \begin{equation}
  \vx = \va - \bar a\,\vone = \va - \vu\,(\vu\T\va) = (\mI - \vu\vu\T)\,\va
  = \mP\,\va.
  \label{eq:contrast}
  \end{equation}
\end{itemize}
In words: \textbf{contrast = each channel minus the pixel's own average}.
Figure~\ref{fig:split} shows it for three real pixels.

$\mP = \mI - \vu\vu\T$ is called a \emph{projector}.  $\vu\vu\T$ is an
$n\times n$ matrix with every entry $1/n$; multiplying a spectrum by it
replaces every channel with the average.  $\mP$ has three properties used
below:
\begin{align}
\mP\vone &= \mathbf 0 && \text{a flat spectrum has no contrast at all;} \label{eq:P1}\\
\mP\mP &= \mP && \text{removing the average twice changes nothing;} \label{eq:PP}\\
\mP\T &= \mP && \text{it is symmetric.} \label{eq:PT}
\end{align}

For all pixels at once, the code computes \code{A - np.outer(A @ u, u)}.
\code{A @ u} is a column of $N$ numbers, $\vu\T\va$ for each pixel; the outer
product with $\vu$ turns each into that pixel's flat part.  So
\begin{equation}
\mX = \mA - (\mA\vu)\vu\T = \mA\,(\mI - \vu\vu\T) = \mA\mP.
\end{equation}
Each row of $\mX$ is one pixel's contrast.  Every row adds up to zero.

\subsection{What it does to the model}

Apply $\mP$ to both sides of the model~\eqref{eq:model}.  By
\eqref{eq:P1} the shading term vanishes:
\begin{equation}
\vx = \mP\va = g\,\underbrace{\mP\vone}_{=\,\mathbf 0} + y\,\mP\epsy +
b\,\mP\epsb = y\,\ey + b\,\eb,
\qquad \ey = \mP\epsy,\ \eb = \mP\epsb.
\label{eq:contrastmodel}
\end{equation}
This is the reason for the step: after it, shading is gone and only the two
dyes remain.

\subsection{What it costs}

$\mP$ also removes each dye's own average.  The contrast
spectra $\ey$, $\eb$ are the dye spectra with their flat part taken out.  A
dye that absorbed equally in every channel would vanish completely: it
cannot be told from a shadow.  The docstring says this plainly (``$y$ and
$b$ measure \emph{contrast}, not total optical density'').
Figure~\ref{fig:split} shows how much it removes from real pixels: the brown
example averages $A = 0.48$, and all of that average is discarded.  The pale
yellow example keeps more of its signal, because yellow absorbs very unevenly
across the channels.  This is why, in \file{pca\_study}, the strength inside
the dot hardly changes with concentration (Section~\ref{sec:limits}).

\begin{figure}[H]
\centering
\includegraphics[width=\linewidth]{duotone-split}
\caption{Three real spectra split by \code{contrast}.  Top: absorbance in
the six channels, with the pixel's average (red dashes).  Bottom: absorbance
minus the average, which is all that later steps see.  \emph{White}: C1034's
pad at least 6~px outside the dot (780 pixels); essentially zero everywhere,
as it should be against the dry pad.  \emph{Yellow}: C1034 (2.2~ppb), at
least 3~px inside the dot (98 pixels); absorbs violet most, then royal blue,
little else.  \emph{Brown}: C1024 (80~ppb), inside the dot (514 pixels);
absorbs strongly from violet to green, not photo red.  Each bar is the median
of the group's pixels.  Made by \file{duotone\_guide.py}.}
\label{fig:split}
\end{figure}

%======================================================================
\section{Step 2: the color map (\code{hue\_plane})}
\label{sec:plane}
%======================================================================

\src{288}{290}{The body of \code{hue\_plane}.}

\subsection{The idea}

After Step~1, each pixel's contrast $\vx$ is a point in six-dimensional space
(five really, since the entries add to zero).  By
equation~\eqref{eq:contrastmodel}, if the model is right every point is
$y\ey + b\eb$: it lies on the flat sheet (plane) through the origin that
contains the two arrows $\ey$ and $\eb$.  \code{hue\_plane} finds that plane
from the data without knowing $\ey$ or $\eb$.  Once the plane is known, every
pixel can be drawn on a two-dimensional map.

\subsection{The algebra: singular value decomposition}

The tool is the singular value decomposition (SVD).  Any $N\times n$ matrix
can be written as
\begin{equation}
\mX = \mU\,\mS\,\mV\T,
\label{eq:svd}
\end{equation}
where the rows of $\mV\T$, called $\vv_1\T,\dots,\vv_n\T$, are $n$
perpendicular unit-length directions in channel space; $\mS$ is diagonal
with $s_1 \ge s_2 \ge \dots \ge 0$; and $\mU$ has perpendicular unit-length
columns.  What each piece means here:
\begin{itemize}
\item \textbf{$\mX\vv_k$} (multiply both sides of \eqref{eq:svd} by
  $\mV$: $\mX\mV = \mU\mS$).  Row $i$ of $\mX\vv_k$ is $\vx_i\T\vv_k$, the
  dot product of pixel $i$ with direction $\vv_k$.  \emph{Meaning:} how far
  pixel $i$ reaches along direction $k$ --- its coordinate on axis $k$.
\item \textbf{$s_k^2$} is the sum over all pixels of that coordinate squared.
  \emph{Meaning:} how much of the whole cloud lies along direction $k$.
\item The directions come out in order: $\vv_1$ captures the most, $\vv_2$
  the most of what is left while being perpendicular to $\vv_1$, and so on.
\end{itemize}
\code{hue\_plane} keeps the first two directions as the rows of a
$2\times n$ matrix $\mB$ (\code{basis}), and reports
\[
\text{explained}_k = \frac{s_k^2}{\sum_j s_j^2},
\]
the share of the cloud along each direction.  Pooled over all 18 pads:
\begin{center}\small
\begin{tabular}{@{}lcccccc@{}}
\toprule
direction $k$ & 1 & 2 & 3 & 4 & 5 & 6 \\
\midrule
share & 89.0\,\% & 9.3\,\% & 0.75\,\% & 0.68\,\% & 0.20\,\% & 0 \\
\bottomrule
\end{tabular}
\end{center}
The first two hold 98.3\,\%, which is what two colors predict.  The sixth
is exactly zero, as a check: every row of $\mX$ is perpendicular to $\vu$
($\mX\vu = \mA\mP\vu = \mathbf 0$), so the cloud has no extent along $\vu$ at
all.  The 1.7\,\% in directions 3--5 is noise plus whatever the two-color
model misses.

\paragraph{Why not ordinary PCA?}  Ordinary PCA first subtracts the
\emph{average pixel} from every pixel, then does the SVD.  It measures spread
around the average.  \code{hue\_plane} skips that subtraction and measures
spread around the \emph{origin}.  Here the origin has a physical meaning: it
is the dry pad (no dye, $A = 0$ in every channel), and the model says every
pixel is a sum of arrows starting there.  The average pixel is a mix of dot
and white and has no such meaning.  So ``explained'' is a share of the total
sum of squares, not of variance in the textbook sense.

\subsection{The map}
\label{sec:map}

A pixel's position on the map is
\begin{equation}
\vs = \mB\,\vx = \begin{pmatrix}\vv_1\T\vx\\ \vv_2\T\vx\end{pmatrix}
\quad\text{and for all pixels}\quad
\mathbf S_2 = \mX\mB\T\ (N\times 2).
\end{equation}
\emph{Meaning of $\mX\mB\T$:} every pixel's dot product with each of the two
map directions, i.e.\ its two map coordinates.  (This is \code{X @ basis.T}
in \code{endmembers} and in \file{pca\_study.py}.)  On the map:
\begin{itemize}[nosep]
\item \textbf{distance} from the origin, $\|\vs\|$, is the \emph{strength}
  of the color;
\item \textbf{angle}, $\operatorname{atan2}(s_2, s_1)$, is the \emph{hue}.
\end{itemize}
This works because of the model: a pixel with twice as much of the same
mixture sits twice as far out along the same line.  The angle depends only on
the ratio $y:b$.

Going the other way, a direction on the map, $\mathbf{q} = (\cos\phi,
\sin\phi)\T$, corresponds to the spectrum
\begin{equation}
\mathbf e = \mB\T\mathbf q = \cos\phi\;\vv_1 + \sin\phi\;\vv_2,
\label{eq:backtocells}
\end{equation}
\emph{meaning:} the contrast spectrum that points in that direction on the
map.  (Code: \code{d @ basis}.)

\begin{figure}[H]
\centering
\includegraphics[width=0.92\linewidth]{duotone-plane}
\caption{All 39{,}789 pad pixels of the 18 pads on the pooled map (gray, log
density).  The axes are the map's axis~1 and axis~2 (the same map as in
\file{pca\_study}; axis~2 is signed so brown is up).  The white crowd at the
origin is the unstained pad.  From it the pixels run out along a straight
line at about $-30^\circ$ (pale yellow) and then bend up to about
$+20^\circ$ (brown).  Red markers: the yellow and brown example pixels of
Figure~\ref{fig:split}.  Arrows: the two reference colors that
\code{duotone.fit} actually returns today; they point away from the data
(Section~\ref{sec:endmembers}).  Inset: the origin enlarged; the dotted
circle is the median strength (0.025), the cutoff \code{endmembers} uses.}
\label{fig:plane}
\end{figure}

%======================================================================
\section{Step 3: the two reference colors (\code{endmembers})}
\label{sec:endmembers}
%======================================================================

\subsection{The idea: the wedge}

By equation~\eqref{eq:contrastmodel} every pixel is $y\ey + b\eb$ with $y,b
\ge 0$.  On the map, that is $y$ times one arrow plus $b$ times another.
Adding two arrows with non-negative amounts can only land \emph{between}
them.  So all pixels lie in a wedge whose two edges point along pure yellow
and pure brown.  The reference colors (``endmembers'') can then be read off
as the two \emph{extreme angles} of the cloud.  Using angles rather than
distances keeps a merely dark pixel from being taken as a pure one.

\subsection{The steps}

\src{318}{340}{The body of \code{endmembers}.}

\begin{enumerate}
\item \textbf{Map coordinates and angles} of every pixel (lines 318--319):
  $\vs_i = \mB\vx_i$, $\theta_i = \operatorname{atan2}(s_{i2}, s_{i1})$.
\item \textbf{Weights} (320--321): by default each pixel's strength
  $w_i = \|\vs_i\|$.
\item \textbf{Middle direction} (325--326).  The average of the pixels'
  unit arrows, weighted by strength:
  $\bar\theta = \operatorname{atan2}\bigl(\sum w_i\sin\theta_i,\ \sum
  w_i\cos\theta_i\bigr)$.  Angles cannot simply be averaged ($+179^\circ$
  and $-179^\circ$ average to $0^\circ$, the wrong way), so the arrows are
  added instead.
\item \textbf{Angles relative to the middle} (327), wrapped into
  $-180^\circ\dots+180^\circ$.  This keeps the wedge from straddling the
  $\pm180^\circ$ seam.
\item \textbf{Keep the stronger half} (328): pixels with $w_i$ above the
  median.  The intent is to leave out the white pixels, whose angle is just
  noise.
\item \textbf{The two edges} (329--330): the 2nd and 98th percentiles of the
  kept pixels' relative angles.  Percentiles instead of minimum and maximum,
  so one stray pixel does not decide.
\item \textbf{Back to channel space} (332--336) by
  equation~\eqref{eq:backtocells}, and scaled to length~1.  These are the
  unit contrast spectra $\ey$ and $\eb$.
\item \textbf{Which is which} (338--340).  An SVD direction's sign is
  arbitrary, so the two edges have to be named.  \code{\_long\_wave\_share}
  shifts each spectrum so its lowest channel is zero, then divides its
  average at 500~nm and above by its average below 500~nm.  Brown keeps
  absorbing into the green where yellow has stopped, so the edge with the
  larger ratio is called brown.
\end{enumerate}

\subsection{What goes wrong today}
\label{sec:endmembersfail}

On the pooled 18 pads the two edges come out at $+150^\circ$ and
$-164^\circ$ on the map (arrows in Figure~\ref{fig:plane}), pointing away
from every stained pixel.  Figure~\ref{fig:spectra} compares them with the
real yellow and brown: each is almost exactly the \emph{negative} of the
real color (the cosine of the angle between them is $-0.995$ for yellow and
$-0.996$ for brown).

\begin{figure}[H]
\centering
\includegraphics[width=\linewidth]{duotone-spectra}
\caption{Red: the unit contrast spectra of the yellow and brown example
pixels of Figure~\ref{fig:split}.  Gray: what \code{duotone.fit} returns as
``yellow'' and ``brown'' on the 18 pads.  Each gray set is nearly a mirror
image of the red: where yellow absorbs most (violet), fit's ``yellow''
absorbs least.}
\label{fig:spectra}
\end{figure}

\paragraph{Why.}  Most of each pad is white: the median strength of all
pixels is only 0.025, about four times the noise (Figure~\ref{fig:plane},
inset).  ``Above the median'' therefore keeps many near-white pixels whose
angle is set by noise and points in every direction; 11\,\% of the kept
pixels point more than $90^\circ$ away from the middle direction.  The 2nd
and 98th percentiles land among them, on the far side of the origin.  The
naming step then names the two wrong-way spectra as best it can.

\paragraph{What it does downstream.}  With both reference colors pointing
backward, \code{unmix} finds almost no yellow or brown anywhere.
Running \file{duotone.py} by itself (its \code{\_\_main\_\_}) today gives:

\begin{center}\small
\begin{tabular}{@{}lrrrrrr@{}}
\toprule
pad & ppb & pixels & white & yellow & brown & residual (A) \\
\midrule
C1034 & 2.2 & 2166 & 2061 & 94 & 11 & 0.054 \\
C1032 & 5 & 2237 & 2125 & 69 & 43 & 0.073 \\
C1017 & 10 & 2166 & 2038 & 32 & 96 & 0.107 \\
C1029 & 20 & 2236 & 2040 & 68 & 128 & 0.081 \\
C1035 & 40 & 2165 & 2032 & 33 & 100 & 0.090 \\
C1024 & 80 & 2261 & 2141 & 46 & 74 & 0.130 \\
C1025 & 160 & 2218 & 2088 & 24 & 106 & 0.158 \\
\bottomrule
\end{tabular}
\end{center}
Between 91 and 96\,\% of pixels are called white on every pad, and what the
model leaves unexplained (0.05--0.16) is 7--21 times the noise ($\approx
0.0075$).
The other 11 pads are similar.  These numbers mean nothing about the pads;
they show the fit has failed.

The same function also runs on each pad alone inside \code{cell\_PCA} (its
default).  \file{pca\_study} step~1 found that it fails there too on pale
pads (e.g.\ C1034: $+152^\circ$ and $+149^\circ$) and works on some dark
ones (C1021: $-36^\circ$ and $+24^\circ$).

\paragraph{How close a fix is.}  Only the ``stronger half'' cutoff is at
fault.  Keeping only pixels with strength above 0.1 instead (13 times the
noise) puts the two edges at
$-40^\circ$ and $+27^\circ$, where the data are; above 0.2, at $-32^\circ$
and $+27^\circ$.  A different route, fixing the two reference colors once
from chosen pixels, is the plan already in
\file{docs/pca\_study\_2026-09-27.md}.

%======================================================================
\section{Step 4: split each pixel (\code{unmix})}
\label{sec:unmix}
%======================================================================

\src{386}{406}{The body of \code{unmix}.}

\subsection{What is being solved}

For each pixel, find the recipe $(g, y, b)$ whose spectrum comes closest to
the measured one, with $y, b \ge 0$:
\begin{equation}
\min_{g,\ y\ge0,\ b\ge0}\ \bigl\|\va - g\vone - y\epsy - b\epsb\bigr\|^2 .
\label{eq:lsq}
\end{equation}
``Closest'' means the smallest sum of squared differences over the channels
(least squares).  Every channel counts equally.

\subsection{Getting rid of $g$}

Split every spectrum into its flat part and its contrast
(Section~\ref{sec:contrast}): $\va = (\vu\T\va)\vu + \vx$ and $\epsy =
(\vu\T\epsy)\vu + \ey$, and the same for brown.  Then the misfit inside
\eqref{eq:lsq} is
\[
\underbrace{\bigl[\vu\T\va - \sqrt n\,g - y\,\vu\T\epsy -
b\,\vu\T\epsb\bigr]\vu}_{\text{flat part}}
\;+\;
\underbrace{\bigl[\vx - y\ey - b\eb\bigr]}_{\text{contrast part}} .
\]
(The $\sqrt n$ appears because $\vone = \sqrt n\,\vu$.)  The two parts are
perpendicular, so the squared length of the sum is the sum of their squared
lengths.  $g$ appears only in the flat part, and whatever $y$ and $b$ are,
$g$ can be chosen to make the flat part exactly zero.  So the problem
becomes:
\begin{equation}
\min_{y\ge0,\ b\ge0}\ \bigl\|\vx - y\ey - b\eb\bigr\|^2
\;=\; \min_{\bm\beta\ge0}\ \|\vx - \mE\bm\beta\|^2,
\qquad
\mE = \begin{pmatrix}\ey & \eb\end{pmatrix}\ (n\times 2),\
\bm\beta = \begin{pmatrix} y\\ b\end{pmatrix}.
\label{eq:lsq2}
\end{equation}
This is lines 386--389 of the code: \code{Xc} is $\mX$, and \code{ey},
\code{eb} are the reference colors with their flat parts removed.

Note what this says: $g$ takes up not only shading but also the flat part of
each dye.  That is the ``cost'' of Section~\ref{sec:contrast} again, now seen
from the fitting side.

\subsection{The best $y$ and $b$, ignoring the ``$\ge 0$''}

At the best $\bm\beta$, the leftover $\vx - \mE\bm\beta$ must be
perpendicular to both reference colors; otherwise a little more of one of
them would fit better.  Written out:
\begin{equation}
\mE\T(\vx - \mE\bm\beta) = \mathbf 0
\quad\Longrightarrow\quad
\underbrace{\mE\T\mE}_{\mM}\,\bm\beta = \underbrace{\mE\T\vx}_{\vr}
\quad\Longrightarrow\quad
\bm\beta = \mM\inv\vr .
\label{eq:normal}
\end{equation}
The meaning of each product:
\begin{itemize}
\item $\vr = \mE\T\vx = \begin{pmatrix}\ey\T\vx\\ \eb\T\vx\end{pmatrix}$
  (\code{rhs}): the dot product of the pixel with each reference.
  \emph{Meaning:} a raw score for how much the pixel resembles yellow, and
  how much it resembles brown.
\item $\mM = \mE\T\mE = \begin{pmatrix}\ey\T\ey & \ey\T\eb\\ \ey\T\eb &
  \eb\T\eb\end{pmatrix}$ (\code{M}): the dot products of the references
  with each other.  \emph{Meaning:} how much the two references resemble
  each other.  The diagonal holds each one's squared length (1 for unit
  references); the off-diagonal is their overlap.
\item $\mM\inv\vr$: \emph{meaning:} the raw scores, untangled.  The two
  references overlap (on these pads their angle is $49^\circ$, overlap
  $\cos 49^\circ = 0.65$), so a brown pixel also scores on yellow, and vice
  versa.  Multiplying by $\mM\inv$ removes that double counting.  If the
  references were perpendicular, $\mM$ would be the identity and the scores
  would already be the amounts.
\end{itemize}
For two unit-length references with overlap $c = \ey\T\eb$, the $2\times2$
inverse gives explicitly:
\begin{equation}
y = \frac{r_y - c\,r_b}{1 - c^2},\qquad
b = \frac{r_b - c\,r_y}{1 - c^2}.
\label{eq:explicit}
\end{equation}
In words: each score, minus the share of it that the other color explains,
scaled back up.

For all pixels at once, the code stacks the scores as rows of an $N\times2$
matrix $\mathbf R = \mX\mE$ and computes $\mathbf R\,(\mM\inv)\T$ (line 393).
Since $\mM$ is symmetric, that is $\mM\inv\vr$ for every row.

\subsection{A worked example}
\label{sec:worked}

Take as references the unit contrast spectra of the yellow and brown example
pixels (Figure~\ref{fig:spectra}, red):
\[
\ey = \begin{pmatrix} 0.748\\ 0.343\\ -0.144\\ -0.289\\ -0.292\\ -0.366
\end{pmatrix},\quad
\eb = \begin{pmatrix} 0.339\\ 0.333\\ 0.177\\ 0.010\\ 0.004\\ -0.862
\end{pmatrix},\quad
c = \ey\T\eb = 0.654 .
\]
Take as the pixel the median of the 319 pixels at least 3~px inside the dot
of C1017 (10~ppb), whose hue on the map is $0^\circ$, between yellow
($-30^\circ$) and brown ($+20^\circ$):
\[
\va = \begin{pmatrix} 0.659\\ 0.616\\ 0.453\\ 0.290\\ 0.295\\ 0.004
\end{pmatrix},\ \bar a = 0.386,\quad
\vx = \va - \bar a\vone = \begin{pmatrix} 0.273\\ 0.230\\ 0.067\\ -0.096\\
-0.092\\ -0.382 \end{pmatrix}.
\]
Scores: $r_y = \ey\T\vx = 0.468$, $r_b = \eb\T\vx = 0.509$.  By
\eqref{eq:explicit}:
\[
y = \frac{0.468 - 0.654 \times 0.509}{1 - 0.654^2} = 0.237,\qquad
b = \frac{0.509 - 0.654 \times 0.468}{1 - 0.654^2} = 0.354 .
\]
The leftover $\vx - y\ey - b\eb$ has an RMS of 0.028 per channel, four
times the single-pixel noise, even though it is the median of 319 pixels
(which should have almost no noise left).  So a 10~ppb dot center is
\emph{not} exactly a mix of these two references; this is the ``bend'' of
the cloud visible in Figure~\ref{fig:plane}.  It is a fact about the pads to
keep in mind, not a fault in the code.

\subsection{Keeping the amounts non-negative}

If $\bm\beta = \mM\inv\vr$ has a negative entry, the pixel lies outside the
wedge between the two references.  The code (lines 395--402) then sets that
amount to zero and refits the other one alone: for example, if $y < 0$,
$y = 0$ and $b = \max(r_b / \eb\T\eb,\ 0)$, the amount of brown that fits
best by itself.  Figure~\ref{fig:unmix} shows the four regions this creates
on the plane of the two references.

\begin{figure}[H]
\centering
\includegraphics[width=0.62\linewidth]{duotone-unmix}
\caption{Schematic of \code{unmix} on the plane of the two reference colors
(directions as on the map: yellow $-30^\circ$, brown $+20^\circ$; lengths
arbitrary).  Inside the wedge (pixel P) the amounts are read off by the
parallelogram rule: go $y$ along yellow, then $b$ along brown.  Outside it,
one amount is set to zero and the pixel is dropped at right angles onto the
other reference (pixel Q).  Beyond the dashed lines, at right angles to the
references, both amounts are zero and the pixel is called white.}
\label{fig:unmix}
\end{figure}

This gives the exact best non-negative answer when the two references are
less than $90^\circ$ apart ($c \ge 0$), as they are on these pads
($49^\circ$).  When they are more than $90^\circ$ apart it can be wrong
(finding~F2, Section~\ref{sec:findings}).

\subsection{What is returned}

\code{unmix} returns $(g, y, b)$ for every pixel.  $y$ and $b$ are in
``contrast absorbance'': with unit-length references, $y$ is the length
along the yellow direction, the same scale as strength on the map.  (The
yellow example pixel itself gives $y = 0.435$; its strength on the map is
0.433.)

The returned $g$ is computed as $\vu\T(\va - y\ey - b\eb)$ (line 405).
Because $\ey$ and $\eb$ have no flat part, this is just $\vu\T\va =
\sqrt n\,\bar a$: $\sqrt6 \approx 2.45$ times the pixel's average absorbance.
It is \emph{not} the $g$ of the model~\eqref{eq:model}, which would be $\bar
a - y\,\bar\varepsilon_y - b\,\bar\varepsilon_b$.  Nothing uses $g$ today
(finding~F3).

\paragraph{\code{residual\_rms}} (lines 409--416) repeats the calculation of
the leftover $\vx - y\ey - b\eb$ for every pixel and returns the root mean
square over all pixels and channels: the typical size of what the two colors
do not explain, to compare with the noise.

%======================================================================
\section{Step 5: label each pixel (\code{classify})}
\label{sec:classify}
%======================================================================

\src{444}{450}{The body of \code{classify}.}

Two cuts, both chosen by the caller:
\begin{itemize}[nosep]
\item \textbf{white} if $y + b <$ \code{floor}.  The caller should use a
  multiple of the noise on $y+b$ (Section~\ref{sec:noise});
  \code{\_\_main\_\_} uses three times it.
\item otherwise \textbf{brown} if the brown fraction $b/(y+b) \ge$
  \code{brown\_at} (default 0.5), else \textbf{yellow}.
\end{itemize}
As the docstring says, yellow and brown are ends of a continuum
(Figure~\ref{fig:plane}), so where to cut is a choice, not a discovery.

%======================================================================
\section{How noisy are the amounts?}
\label{sec:noise}
%======================================================================

Here ``noise'' means the pixel scatter of Section~\ref{sec:padnoise}, which
is mostly the paper's own texture.

\subsection{\code{amount\_sigma}: carrying channel noise into $y$ and $b$}

\src{469}{479}{The body of \code{amount\_sigma}.}

Inside the wedge, \code{unmix} is a fixed linear recipe.  From
\eqref{eq:contrast} and \eqref{eq:normal}:
\begin{equation}
\bm\beta = \mM\inv\mE\T\vx = \underbrace{\mM\inv\mE\T}_{\mL}\,\mP\,\va .
\end{equation}
\emph{Meaning of $\mL$} (\code{pinv}, a $2\times n$ matrix): its first row is
a set of weights, one per channel, whose weighted sum of a pixel's contrast
is $y$; the second row does the same for $b$.  For the worked example's
references:
\[
\mL = \begin{pmatrix}
0.919 & 0.219 & -0.452 & -0.516 & -0.514 & 0.344\\
-0.261 & 0.189 & 0.473 & 0.347 & 0.340 & -1.087
\end{pmatrix}.
\]
Each row adds up to zero, so adding the same number to every channel
(shading) cannot change $y$ or $b$.

If the pixel's absorbances have random errors with covariance matrix $\mC$
($n\times n$: channel variances on the diagonal, and off the diagonal how
much two channels' errors tend to move together), then any fixed linear
recipe $\mathbf W$ applied to them produces errors with covariance
$\mathbf W\mC\mathbf W\T$.  Here $\mathbf W = \mL\mP$, so
\begin{equation}
\operatorname{Cov}\begin{pmatrix}y\\b\end{pmatrix}
= \mL\,(\mP\mC\mP\T)\,\mL\T .
\label{eq:sandwich}
\end{equation}
Meaning of the pieces: $\mP\mC\mP\T$ is the noise that survives the removal
of the flat part (the code's \code{C}); the outer $\mL \dots \mL\T$ carries
it onto the two amounts.  The result is $2\times2$: the diagonal holds the
variances of $y$ and $b$, and the off-diagonal how their errors move
together.  For the total,
\[
\sigma_{y+b}^2 = \operatorname{var} y + \operatorname{var} b +
2\operatorname{cov}(y, b),
\]
which is line 478.  (The $\mP$ is in fact redundant, since the rows of $\mL$
already add to zero, so $\mL\mP = \mL$; it does no harm.)

\code{amount\_sigma} uses a \emph{diagonal} $\mC$, $\operatorname{diag}
(\sigma_1^2, \dots, \sigma_n^2)$: it assumes the channels' errors are
unrelated.  With the worked example's references and the pooled $\sigma$s
it gives $\sigma_y = 0.0086$, $\sigma_b = 0.0096$, $\sigma_{y+b} = 0.0076$.

\subsection{\code{noise\_cov}: do the channels' errors move together?}

\src{506}{512}{The body of \code{noise\_cov}.}

This builds the full $\mC$.  The correlation $\rho_{jk}$ between channels $j$
and $k$ (a number from $-1$ to $1$: 0 means unrelated, 1 means always moving
together) is measured from the neighbor differences \code{dnoise}; the
sizes $\sigma_j$ from the robust per-channel values.  Then
\[
C_{jk} = \sigma_j\,\sigma_k\,\rho_{jk}
\qquad\text{(code: \code{np.outer(sigma, sigma) * corr})}.
\]
Pooled over the 18 dry pads:
\begin{center}\small
\begin{tabular}{@{}lrrrrrr@{}}
\toprule
 & violet & r.\ blue & cyan & green Gb & green Gr & p.\ red \\
\midrule
violet    & 1 & 0.20 & 0.17 & 0.01 & 0.11 & 0.02 \\
royal blue&   & 1    & 0.33 & 0.06 & 0.16 & 0.03 \\
cyan      &   &      & 1    & 0.13 & 0.22 & 0.05 \\
green Gb  &   &      &      & 1    & 0.11 & 0.16 \\
green Gr  &   &      &      &      & 1    & 0.12 \\
photo red &   &      &      &      &      & 1    \\
\midrule
$\sigma$ (A) & 0.0062 & 0.0066 & 0.0073 & 0.0064 & 0.0064 & 0.0075 \\
\bottomrule
\end{tabular}
\end{center}
The correlations are weak (0.01 to 0.33).  With this full $\mC$ in
\eqref{eq:sandwich}, $\sigma_{y+b}$ becomes 0.0079 instead of 0.0076: the
diagonal shortcut makes little difference here.

The \code{noise\_cov} docstring used to say that most pixel-to-pixel
variation dims every channel together, so the correlations are ``the
point''.  That was measured on the bare paper.  On the dry pad the answer
depends on the scale.  Neighbor differences, as used here, correlate only
weakly between channels (the table above).  But texture 1--4~px across
agrees between different LEDs at 0.5--0.8 (shading study, step~1a).  The
docstring now says both (reworded 2026-09-29; finding~F4).

%======================================================================
\section{Choosing channels: \code{j\_hue}, \code{j\_amounts},
\code{subset\_sigma}, \code{subset\_unmix}}
\label{sec:subsets}
%======================================================================

These four answer: if only some of the channels were measured, how well
could yellow and brown still be told apart and measured?  Today only
\file{report14.py} (a record that no longer runs) calls them.  They are the
natural tools for step~4 of the suggested order in
\file{docs/pca\_study\_2026-09-27.md} (rank channel subsets and choose the
channels), which is why they are explained in full.  In all four, \code{idx} picks the channels and
$\mC$ is the full noise covariance, restricted to those channels.  Below,
$\vone$ is all ones on the chosen channels.

\subsection{Weighted least squares}

\code{unmix} counts every channel equally.  These functions instead weight
by the noise: they minimize
\begin{equation}
(\va - \mF\bm\theta)\T\,\mC\inv\,(\va - \mF\bm\theta).
\label{eq:wls}
\end{equation}
\emph{Meaning of the $\mC\inv$ in the middle:} a channel with twice the noise
counts a quarter as much, and two channels whose errors move together are not
counted as two independent witnesses.  With plain least squares, $\mC\inv$
is replaced by the identity.  The same argument as \eqref{eq:normal} gives
\begin{equation}
\hat{\bm\theta} = (\mF\T\mC\inv\mF)\inv\,\mF\T\mC\inv\va,
\qquad
\operatorname{Cov}(\hat{\bm\theta}) = (\mF\T\mC\inv\mF)\inv .
\label{eq:wlssol}
\end{equation}
$\mF\T\mC\inv\mF$ is called the \emph{information matrix}.  \emph{Meaning:}
how sharply the data pin down each amount.  More information means smaller
errors; its inverse is the error covariance of the estimates.

\subsection{\code{j\_hue}: can the channels tell yellow from brown?}

\src{541}{549}{The body of \code{j\_hue}.}

Take any weighted sum of the channels, $t = \vw\T\va$, as an index.  Ask
three things of it:
\begin{itemize}[nosep]
\item \textbf{Shading-proof:} adding the same amount to every channel must
  not change it: $\vw\T\vone = 0$.
\item \textbf{Signal:} how much it differs between a unit amount of pure
  brown and a unit amount of pure yellow: $\vw\T(\epsb - \epsy) = \vw\T\vd$,
  with $\vd = \epsb - \epsy$.
\item \textbf{Noise:} its random error, $\sqrt{\vw\T\mC\vw}$ (the sandwich
  rule again, for one row).
\end{itemize}
$J$ is the best achievable signal-to-noise ratio:
\[
J = \max_{\vw:\ \vw\T\vone = 0}\ \frac{|\vw\T\vd|}{\sqrt{\vw\T\mC\vw}} .
\]
\textbf{Derivation.}  Scaling $\vw$ does not change the ratio, so fix the
noise at $\vw\T\mC\vw = 1$ and maximize $\vw\T\vd$ with $\vw\T\vone = 0$.
Using Lagrange multipliers $\mu$ and $\lambda$ for the two conditions, the
best $\vw$ satisfies
\[
\vd - \mu\,\mC\vw - \lambda\vone = \mathbf 0
\quad\Longrightarrow\quad
\vw \propto \mC\inv(\vd - \lambda\vone) .
\]
The shading condition $\vone\T\vw = 0$ fixes $\lambda$:
\[
\vone\T\mC\inv\vd - \lambda\,\vone\T\mC\inv\vone = 0
\quad\Longrightarrow\quad
\lambda = \frac{\vone\T\mC\inv\vd}{\vone\T\mC\inv\vone} .
\]
Let $\vz = \vd - \lambda\vone$, so $\vw \propto \mC\inv\vz$ and $\vone\T
\mC\inv\vz = 0$.  Then $\vw\T\vd \propto \vz\T\mC\inv\vd = \vz\T\mC\inv(\vz +
\lambda\vone) = \vz\T\mC\inv\vz$, and $\vw\T\mC\vw \propto \vz\T\mC\inv\vz$
also.  So
\begin{equation}
J^2 = \frac{(\vz\T\mC\inv\vz)^2}{\vz\T\mC\inv\vz} = \vz\T\mC\inv\vz
= \vd\T\mC\inv\vd - \frac{(\vone\T\mC\inv\vd)^2}{\vone\T\mC\inv\vone} .
\label{eq:jhue}
\end{equation}
This is lines 546--549 (\code{a - b*b/c}).  \emph{Meaning of the two terms:}
$\vd\T\mC\inv\vd$ is how many noise-sizes separate brown from yellow if
shading could be ignored (the difference measured in units of noise).  The
subtracted term is the part of that separation that lies along the shading
direction, as seen through the noise, and is given up to stay shading-proof.
With a single channel the two terms are equal and $J = 0$: one channel alone
cannot tell a lot of yellow from a little brown.

With all six channels and the example references, $J = 130$ per unit amount;
at a typical dot strength of 0.5, pure yellow and pure brown differ by about
65 noise-sizes in a single pixel.

\subsection{\code{j\_amounts} and \code{subset\_sigma}: how precisely is
each amount known?}

\src{575}{587}{The body of \code{j\_amounts}.}

Here the full model is used, $\mF = (\epsy\ \epsb\ \vone)$, with amounts
$(y, b, g)$.  The information matrix \eqref{eq:wlssol} splits into blocks,
with $\mE = (\epsy\ \epsb)$:
\[
\mF\T\mC\inv\mF =
\begin{pmatrix}
\mE\T\mC\inv\mE & \mE\T\mC\inv\vone\\
\vone\T\mC\inv\mE & \vone\T\mC\inv\vone
\end{pmatrix}
\quad
\begin{matrix} \leftarrow\ (y, b)\ \text{rows}\\ \leftarrow\ g\ \text{row}
\end{matrix}
\]
We want only the error covariance of $(y, b)$, i.e.\ the top-left
$2\times2$ block of the \emph{inverse}.  The standard formula for the inverse
of a block matrix gives it as the inverse of
\begin{equation}
\mM = \mE\T\mC\inv\mE \;-\;
\frac{(\mE\T\mC\inv\vone)(\vone\T\mC\inv\mE)}{\vone\T\mC\inv\vone}
\label{eq:schur}
\end{equation}
(called the \emph{Schur complement}), which is line 583.  \emph{Meaning:}
the information about $(y,b)$, minus the part of it that has to be spent on
pinning down the unknown shading $g$.  It has the same shape as
\eqref{eq:jhue} for the same reason.  Then $\mM\inv$ is the $2\times2$ error
covariance of $(y,b)$; \code{j\_amounts} returns $1/\sigma_y$ and
$1/\sigma_b$, and \code{subset\_sigma} returns $\sigma_y$, $\sigma_b$ and
$\sigma_{y+b}$.  With all six channels and the example references,
\code{subset\_sigma} gives $\sigma_y = 0.0082$, $\sigma_b = 0.0092$,
$\sigma_{y+b} = 0.0078$, close to \code{amount\_sigma}'s.

\paragraph{Why at least three channels.}  Removing the shading uses up one
channel's worth of information (Section~\ref{sec:contrast}), and two amounts
need two more.  With two channels $\mM$ cannot be inverted; the functions
return zero precision (\code{j\_amounts}) or infinite error
(\code{subset\_sigma}).  \code{subset\_sigma(..., shading=False)} drops the
subtracted term of \eqref{eq:schur}, which assumes the shading is known,
so that two channels can be priced.

\subsection{\code{subset\_unmix}}

\src{648}{663}{The body of \code{subset\_unmix}.}

This applies \eqref{eq:wlssol} to every pixel at once, using only the chosen
channels.  \code{np.linalg.pinv(M) @ (E.T @ Ci @ X.T)} is
$(\mF\T\mC\inv\mF)\inv\mF\T\mC\inv$ applied to every pixel's spectrum (the
columns of $\mX\T$).  Unlike \code{unmix}, it does not remove the flat part
first; it fits $g$ as a third amount, which comes to the same thing.  It then
sets negative amounts to zero \emph{without} refitting the other amount, so
it is not the same as \code{unmix}'s clamp (finding~F5).

%======================================================================
\section{The other pieces}
\label{sec:helpers}
%======================================================================

\begin{description}
\item[\code{build(names)}] makes one \code{PadCube} per pad and keeps only
  the channels every pad has.  A matrix with a hole in it is worse than a
  narrower one.
\item[\code{fit(cubes)}] stacks the contrast of every pad's pixels into one
  tall $\mX$, then runs \code{hue\_plane} and \code{endmembers} on it, so
  that ``brown'' means the same spectrum on every pad.  (Today it inherits
  the failure of Section~\ref{sec:endmembersfail}.)
\item[\code{render\_rgb}] makes a picture for people to look at.  For each
  camera color layer it picks the brightest white-LED exposure that does not
  clip on the pad or the gray chip, divides by the gray chip's mean (so the
  gray chip comes out neutral), stacks red, green (the average of the two
  green layers) and blue, scales so the 99.5th percentile is full white, and
  applies a display gamma.  It plays no part in any measurement.
\item[\code{cell\_name}, \code{cell\_nm}, \code{NM}] spell a channel as
  \code{"led.layer"} and give the LED's nominal wavelength.
\item[\code{HUE\_SPLIT\_NM} = 500] the dividing line used to name the two
  reference colors (Section~\ref{sec:endmembers}, step 8).  With the six
  channels, ``below'' is violet and royal blue; ``500 and above'' is cyan,
  both greens and photo red.
\item[\code{SATURATED\_DN} = 1000] the clipping level for
  \code{render\_rgb} only.
\item[\code{\_\_main\_\_}] builds all 18 pads, fits one pair of reference
  colors, and prints each pad's white/yellow/brown counts and residual
  (the table in Section~\ref{sec:endmembersfail}).
\end{description}

%======================================================================
\section{Limits of the model}
\label{sec:limits}
%======================================================================

These are properties of the approach, not bugs.

\begin{enumerate}
\item \textbf{A color that absorbs evenly is invisible.}  Removing the flat
  part removes shading, and with it the even part of every dye.  Brown
  absorbs fairly evenly from violet to green (Figure~\ref{fig:split}), so
  much of its darkness is discarded; what survives is mostly its drop at
  photo red.
  That is why, in \file{pca\_study} step~2b, the strength inside the dot is
  0.44--0.59 on every pad from 2.2 to 160~ppb, and concentration shows up as
  a change of hue rather than of strength.  For estimating arsenic, the
  discarded flat part may carry information that is worth recovering some
  other way (for example, shading measured on the dry pad in the pre-image
  and subtracted, rather than fitted).  This has since been tried
  (\file{docs/shading\_study\_2026-09-29.md}): the pre-image removes about two
  thirds of the rim's shadow in each channel, and about four fifths on the
  six-channel average, but only 10\,\% of the pixel-scale texture inside the
  pad.  Whether the dot's strength then follows concentration, once the flat
  part is no longer removed, is that study's step~2, not yet done.
\item \textbf{Absorbances add only approximately in paper.}  Light
  scattered back out of paper does not obey Beer--Lambert exactly.
  Kubelka--Munk theory says $K/S = (1 - R_\infty)^2 / 2R_\infty$ is the
  quantity that adds, and \file{padphot.py} already computes it.  At the
  strongest dots ($A \approx 0.66$, $R_\infty \approx 0.22$) the difference
  is not small.  Whether $A$ or $K/S$ makes the two colors add better can be
  tested on these pads.
\item \textbf{Two colors, fixed spectra.}  The model has exactly two
  colors, the same on every pixel.  The residual of the worked example
  (Section~\ref{sec:worked}) and the bend in Figure~\ref{fig:plane} show the
  real pads are not exactly that; the chemistry allows at least two, perhaps
  three, products.
\item \textbf{Six channels, five usable directions.}  Removing the flat part
  leaves five independent directions.  Two colors need two of them; the other
  three hold noise and model misfit, and are what \code{residual\_rms}
  measures.
\end{enumerate}

%======================================================================
\section{Review findings}
\label{sec:findings}
%======================================================================

Most serious first.  None has been changed; each is for you to decide.

\begin{enumerate}[label=\textbf{F\arabic*}, leftmargin=2.4em]
\item \textbf{\code{endmembers} returns reversed colors} (lines 328--330).
  The ``stronger half'' of the pixels is mostly noise, so the edges land
  behind the origin (Section~\ref{sec:endmembersfail}).  This breaks
  \code{fit}, \code{\_\_main\_\_}, and \code{cell\_PCA}'s default per-pad
  fit on pale pads.  \emph{Proposal:} either replace the median cutoff by a
  strength cutoff tied to the noise (e.g.\ 10--20 times $\sigma$; 0.1 gives
  $-40^\circ$ and $+27^\circ$), or stop fitting and use fixed reference
  colors (the plan in \file{pca\_study}).  Also, \code{endmembers} could
  check that the two edges are less than $180^\circ$ apart and lie on the
  data's side of the origin, and refuse otherwise.
\item \textbf{\code{unmix}'s clamp is exact only when the references are less
  than $90^\circ$ apart} (lines 397--402).  If both unconstrained amounts
  are negative and the references are more than $90^\circ$ apart, the true
  best answer may be ``yellow only'', but the code returns zero.  On random
  test data with references $128^\circ$ apart, 198 of 2000 pixels disagreed
  with a standard non-negative solver; with the real references ($49^\circ$
  apart), none did.  \emph{Proposal:} evaluate all candidates (both, yellow
  only, brown only, none) and keep the best, or assert $c \ge 0$.
\item \textbf{The returned $g$ is not the model's $g$} (line 405).  It is
  $\sqrt n$ times the pixel's average absorbance, including the dyes' flat
  parts.  Nothing uses it.  \emph{Proposal:} correct the docstring to say
  what it is, or return $\bar a$ and call it that.
\item \textbf{Noise statements out of date --- reworded 2026-09-29.}
  \code{amount\_sigma}'s comment on line 471 mixed up ``between pixels'' and
  ``between channels''; it now says the channels are taken as unrelated.
  \code{noise\_cov}'s docstring said most variation is common to all
  channels, which was measured on the bare paper; it now gives both dry-pad
  measurements (Section~\ref{sec:noise}).  \code{PadCube.sigma} and
  \code{\_pair\_noise} now call $\sigma$ what it is, mostly the paper's
  texture (Section~\ref{sec:scatterwhat}).  What remains:
  \code{amount\_sigma} still uses the diagonal covariance, a small effect
  (0.0076 against 0.0079).  \emph{Proposal:} optionally let it take the full
  covariance.
\item \textbf{\code{subset\_unmix}'s clamp does not refit} (lines 659--663),
  unlike \code{unmix}.  Only \file{report14} uses it.  \emph{Proposal:} fix
  when channel selection (\file{pca\_study} step~4) needs it.
\item \textbf{Small documentation slips.}
  \code{subset\_sigma}'s docstring mentions ``all 26 cells'' (now 6) and has
  a stray ``Args (continued):'' heading.  \file{analysis\_1.py} line 800 and
  \file{padphot.py} line 57 point to \code{duotone.PadCube.\_reflectance},
  which no longer exists.  \code{hue\_plane} calls its output a
  ``variance fraction''; it is a share of the sum of squares about the
  origin (Section~\ref{sec:plane}).
\item \textbf{\code{unmix} weights all channels equally}, while the
  \code{subset\_*} functions weight by noise.  With six channels of similar
  noise (0.0062--0.0075) this hardly matters; it would matter if channels
  with very different noise were added.  No action needed now.
\end{enumerate}

%======================================================================
\appendix
\section{Symbols}
%======================================================================

\begin{center}\small
\begin{tabular}{@{}lL{11cm}@{}}
\toprule
$n$, $N$ & number of channels (6); number of pad pixels \\
$\va$ & one pixel's absorbances, a column of $n$ \\
$\mA$ & all pixels, $N\times n$, one row per pixel (\code{PadCube.A}) \\
$\vone$, $\vu$ & all ones; the same scaled to length 1 (\code{flat\_direction}) \\
$\bar a$ & a pixel's average over channels \\
$\mP = \mI - \vu\vu\T$ & removes a spectrum's average (projector) \\
$\vx = \mP\va$, $\mX = \mA\mP$ & contrast of one pixel; of all pixels (\code{contrast}) \\
$g, y, b$ & amounts of shading, yellow, brown \\
$\epsy, \epsb$ & yellow and brown spectra \\
$\ey, \eb$ & the same with their average removed; unit length when from \code{endmembers} \\
$\mE = (\ey\ \eb)$ & the two references side by side, $n\times2$ \\
$\mM = \mE\T\mE$ & how much the references overlap, $2\times2$ \\
$\vr = \mE\T\vx$ & a pixel's raw resemblance to each reference \\
$c = \ey\T\eb$ & overlap of unit references (cosine of the angle between them) \\
$\mL = \mM\inv\mE\T$ & weights turning contrast into amounts (\code{pinv}) \\
$\mB$ & the two map directions, $2\times n$ (\code{basis}) \\
$\vs = \mB\vx$ & a pixel's position on the map \\
$\sigma_k$ & pixel scatter of channel $k$ (mostly the paper's texture) \\
$\mC$ & noise covariance between channels, $n\times n$ (\code{noise\_cov}) \\
$\mF = (\vone\ \epsy\ \epsb)$ & all three ingredients, $n\times3$ \\
$J$ & best signal-to-noise for telling brown from yellow (\code{j\_hue}) \\
\bottomrule
\end{tabular}
\end{center}

\paragraph{Reproducing the figures and numbers.}  The figures come from
\file{desktop/padanalysis/duotone\_guide.py} (\code{run()}), which writes
\file{tex/figs/duotone-*.pdf}.  The table in
Section~\ref{sec:endmembersfail} is the output of running
\file{duotone.py} by itself.  Line numbers in the listings are those of
\file{duotone.py} on 2026-09-28.

\end{document}
